% Options for packages loaded elsewhere
\PassOptionsToPackage{unicode}{hyperref}
\PassOptionsToPackage{hyphens}{url}
\PassOptionsToPackage{dvipsnames,svgnames,x11names}{xcolor}
\documentclass[
  11pt,
]{article}
\usepackage{xcolor}
\usepackage[margin=27mm]{geometry}
\usepackage{amsmath,amssymb}
\setcounter{secnumdepth}{-\maxdimen} % remove section numbering
\usepackage{iftex}
\ifPDFTeX
  \usepackage[T1]{fontenc}
  \usepackage[utf8]{inputenc}
  \usepackage{textcomp} % provide euro and other symbols
\else % if luatex or xetex
  \usepackage{unicode-math} % this also loads fontspec
  \defaultfontfeatures{Scale=MatchLowercase}
  \defaultfontfeatures[\rmfamily]{Ligatures=TeX,Scale=1}
\fi
\usepackage{lmodern}
\ifPDFTeX\else
  % xetex/luatex font selection
  \setmainfont[]{Cambria}
  \setmonofont[]{Consolas}
\fi
% Use upquote if available, for straight quotes in verbatim environments
\IfFileExists{upquote.sty}{\usepackage{upquote}}{}
\IfFileExists{microtype.sty}{% use microtype if available
  \usepackage[]{microtype}
  \UseMicrotypeSet[protrusion]{basicmath} % disable protrusion for tt fonts
}{}
\makeatletter
\@ifundefined{KOMAClassName}{% if non-KOMA class
  \IfFileExists{parskip.sty}{%
    \usepackage{parskip}
  }{% else
    \setlength{\parindent}{0pt}
    \setlength{\parskip}{6pt plus 2pt minus 1pt}}
}{% if KOMA class
  \KOMAoptions{parskip=half}}
\makeatother
\setlength{\emergencystretch}{3em} % prevent overfull lines
\providecommand{\tightlist}{%
  \setlength{\itemsep}{0pt}\setlength{\parskip}{0pt}}
\usepackage{bookmark}
\IfFileExists{xurl.sty}{\usepackage{xurl}}{} % add URL line breaks if available
\urlstyle{same}
\hypersetup{
  pdftitle={Sequential Composition of Leakage-Budgeted Agents: an Exponential-to-Linear Separation under Structural Context Erasure},
  pdfauthor={The Privacy-is-Value Research Programme},
  colorlinks=true,
  linkcolor={black},
  filecolor={Maroon},
  citecolor={Blue},
  urlcolor={blue},
  pdfcreator={LaTeX via pandoc}}

\title{Sequential Composition of Leakage-Budgeted Agents: an
Exponential-to-Linear Separation under Structural Context Erasure}
\usepackage{etoolbox}
\makeatletter
\providecommand{\subtitle}[1]{% add subtitle to \maketitle
  \apptocmd{\@title}{\par {\large #1 \par}}{}{}
}
\makeatother
\subtitle{WP-07 · tier A · 2026-07-07 · status revision-draft-v3}
\author{The Privacy-is-Value Research Programme}
\date{2026-07-07}

\begin{document}
\maketitle

\subsection{0. Abstract}\label{abstract}

Multi-agent systems built from sequentially invoked language-model
agents leak information about the sensitive contexts they process, and
system-scale measurement locates the dominant share of that leakage in
the channels between agents rather than in final outputs. We give an
information-theoretic account of when per-agent leakage budgets compose
and when they do not. In a model of N sequentially invoked agents with
declared interfaces and explicit carried state, the best guarantee
available under the budget discipline the empirical literature measures
(marginal, output-only constraints) is geometric: the known (2\^{}N −
1)ε bound for the adversary that observes only the final output is an
upper bound, not an attained leakage, and we exhibit a two-agent
composition in which every marginal budget is zero yet the adversary
that observes the whole transcript recovers a context completely. We
then prove a conditional linear cap: under five numbered engineering
requirements (declared channels only, structural context erasure at
every invocation boundary, fresh local randomness, context provenance,
and conditional interface-inclusive budgets ε\_1, \ldots, ε\_N),
transcript leakage about a source X satisfies I(T; X) ≤ Σ\_i ε\_i. Three
necessity constructions show that the erasure requirement, the
conditional budget form, and the provenance requirement are individually
indispensable, and the cap is achieved, so it is tight. The
exponential-to-linear comparison is a comparison of upper-bound
guarantees under their stated budget semantics, not of attainable
leakages, and the conditional budget at the same numerical level is the
more demanding constraint; the unconditional separation is the
zero-budget construction. The central result is a decomposition, not an
unconditional ceiling. The structural requirements license additive
budget accounting and a Fano error floor. A strict sub-unity bound on
reconstruction additionally requires the declared capacity-deficit
condition Σ\_i ε\_i(t) \textless{} H(X): a numerical fact about a
specific system-adversary pair at a stated time, not a consequence of
the architecture, and the component of any deployed guarantee that
expires as adversary capability grows. Separately, against an informed
adversary holding background side information B\_t that accumulates with
calendar time under the Markov chain B\_t → X → T, the cap survives
conditioning while the error floor erodes through the residual entropy
H(X \textbar{} B\_t): the strict component's deficit condition becomes
Σ\_i ε\_i \textless{} H(X \textbar{} B\_t), a second, genuinely
time-varying expiry channel that requires no growth in what the
transcript reveals (Corollary 5.4b).

\subsection{1. Introduction}\label{introduction}

Language-model agents are now deployed in pipelines: one agent's output
becomes working material for the next, across invocations that may be
separately built, separately budgeted, and separately audited.
Measurement at system scale shows where such pipelines leak. In a recent
benchmark across production LLMs, the channels between agents were
observable in deployed frameworks and carried the dominant share of
leakage about sensitive contexts, and audits that inspect only final
outputs missed a substantial fraction of violations
{[}elyagoubi2026agentleak{]} (exact figures are recorded at Definition
3.6).

The theory matches the measurement. Asif and Amiri prove that in a
sequential pipeline satisfying their Markov structure, with mutually
independent sensitive variables and a marginal per-agent budget I(O\_i;
S\_i) ≤ ε\_i on each agent, leakage about the joint secrets at the final
output is at most the geometric sum Σ\_i 2\^{}\{N−i\} ε\_i, that is
(2\^{}N − 1)ε in the uniform case; that geometric sum is the guarantee
the marginal discipline provides, an upper bound with no matching
attainability construction known, and their measurements show mutual
information rising with chain depth {[}asif2026infotheoretic{]};
independent composition analysis reaches the same qualitative conclusion
{[}patil2025composition{]}. Per-agent budget certification, in the form
the current literature states it, does not compose: N certified agents
are not a certified system.

This paper asks what discipline restores linear accounting, and against
which adversary. Two distinctions organise the answer. The first is the
adversary. Bounds stated against the final output do not survive the
passage to the adversary that observes the whole transcript, and the
transcript adversary is the one the measurements say exists. Proposition
4.3 exhibits a two-agent composition satisfying structural context
erasure and zero marginal budgets, including interface-inclusive
marginal budgets, in which the transcript adversary nevertheless
recovers a context exactly; it is the formal counterpart of the measured
dominance of inter-agent-channel leakage. The second is the form of the
budget: what it ranges over (the public output alone, or everything the
invocation emits including its forward interface) and whether it is
conditioned on what the invocation received. Marginal budgets fail even
at budget zero; conditioning on the received interface is what prices
the correlation the marginal misses (Section 3.4).

The main result is a conditional linear cap (Theorem 5.1): under five
numbered engineering requirements (ER-1 declared channels only, ER-2
structural context erasure, ER-3 fresh local randomness, ER-4 context
provenance, ER-5 conditional interface-inclusive budgets), transcript
leakage about the source is at most the sum of the declared budgets, Nε
in the uniform case. Three of the five hypotheses are individually
load-bearing, each with a necessity construction defeating the cap at
declared budget zero. Dropping the conditional budget form while keeping
erasure caps nothing (Proposition 4.3); dropping erasure while keeping
zero budgets caps nothing (Proposition 5.5, a persistent-pad
construction with total leakage at zero declared budget); dropping
context provenance while keeping erasure and zero budgets caps nothing
(Proposition 5.5b, in which the environment routes the source around the
priced interface through context assembly). No necessity construction is
offered for ER-1 or ER-3. By Theorem 5.1 the five are jointly
sufficient, and Proposition 5.6 shows the cap is achieved. The
exponential-to-linear separation is therefore a separation between
disciplines, not mechanisms, and a comparison of upper-bound guarantees
under stated budget semantics, not of attainable leakages (Remark 5.7).

Throughout, the results are stated inside a decomposition discipline
(Remark 3.4). The structural requirements license two things: additive
budget accounting, and a Fano error floor for any estimator of the
source from the transcript (Corollary 5.4). They do not by themselves
place the leakage-to-entropy ratio below one. The strict claim that
reconstruction is incomplete requires, in addition, the declared
capacity-deficit condition Σ\_i ε\_i(t) \textless{} H(X), evaluated
against a stated decoder class at a stated time t. That condition is a
measurable numerical fact about a system-adversary pair, not a
consequence of the architecture; it is the component that expires as
decoder classes grow, and the paper never states the strict bound
without it.

Our contributions are:

\begin{enumerate}
\def\labelenumi{\arabic{enumi}.}
\item
  A composition model for sequentially invoked agents in which carried
  state is explicit and structural context erasure is a property of the
  composition operator rather than a notational assumption (Definitions
  3.5 to 3.8), together with a taxonomy of budget forms (marginal versus
  conditional, output-only versus interface-inclusive) that the results
  show to be the load-bearing distinctions (Section 3.4).
\item
  An account of the exponential regime: the amplification bound imported
  with its hypotheses carried exactly (Theorem 4.1
  {[}asif2026infotheoretic{]}), an in-model reproduction under the named
  local hypotheses of no carried state, independent noise and bounded
  contamination, the last formally incomparable to the source's
  hypotheses (Proposition 4.2), and the observation that marginal
  budgets do not bind the transcript adversary even at budget zero
  (Proposition 4.3).
\item
  A conditional linear cap (Theorem 5.1): transcript leakage at most
  Σ\_i ε\_i under ER-1 to ER-5, with necessity constructions for three
  of the five hypotheses, the budget form, erasure and provenance
  (Propositions 4.3, 5.5 and 5.5b), tightness (Proposition 5.6), and the
  composed error floor with its declared, time-indexed deficit condition
  (Corollary 5.4).
\item
  A framing discipline that separates what the architecture guarantees
  (additivity and the floor) from what must be declared and expires (the
  deficit), with the time-indexed budgets of ER-6 given precedent in the
  fading-wiretap literature {[}gopala2008secrecy{]}, together with a
  side-information corollary (Corollary 5.4b) separating ER-6's
  certification clock, whose drift is epistemic, from the informed
  adversary's erosion clock, the one time-varying quantity in the paper
  whose drift is not a statement about certification.
\item
  A formally stated open problem (Open Problem 5.8) on the criterion
  separating recoverable from unrecoverable erasure, carried with its
  proof obligations and its unstarted first step named.
\end{enumerate}

Every result in this paper is conditional on the numbered requirements
of Section 3.5, which are obligations on implementations, not facts
about deployed systems; measurement indicates deployed frameworks
routinely violate the first of them today {[}elyagoubi2026agentleak{]}.
Whether a given system satisfies ER-1 to ER-5 is an empirical
conformance question, testable per instance and answered here for no
system. Section 6 states the threats to validity, including the
commitment that measured results against these predictions are reported
with the same prominence as confirmations.

The paper is organised as follows. Section 2 positions the results
against the wiretap, converse, multi-agent leakage, covert-channel,
information-flow, differential-privacy and time-varying-capacity
literatures. Section 3 gives the model, the budget taxonomy, and the
engineering requirements. Section 4 treats the regime without erasure,
including the imported amplification bound and the transcript-adversary
construction. Section 5 proves the linear cap, its necessity and
tightness results, and states the one open problem. Section 6 records
threats to validity, and Section 7 concludes.

\subsection{2. Related work}\label{related-work}

\textbf{The wiretap lineage.} The objects this paper computes with are
the classical ones. Wyner introduced the wiretap channel and
equivocation at the adversary as the measure of secrecy
{[}wyner1975wiretap{]}; Leung-Yan-Cheong and Hellman characterised the
secrecy capacity of the Gaussian channel as a difference of capacities
{[}leungyancheong1978gaussian{]}; Csiszár and Körner treated broadcast
channels with confidential messages, the general single-eavesdropper
form of the family {[}csiszar1978broadcast{]}. Section 3.2 states the
two-agent base results as in-model instances of this family rather than
as new results. The failure of additive accounting for colluding
observers, by contrast, is elementary and needs no citation: joint
information can exceed the sum of marginal informations, with the
interaction identity in Theorem 3.1 locating the excess and Proposition
5.5 exhibiting the extremal case at zero marginals, and requirement ER-1
exists to exclude exactly this failure mode.

\textbf{Fano and the source-coding converses.} The error floor is Fano's
inequality {[}fano1961transmission{]} in its standard converse form
{[}cover2006elements{]}, applied at the level of the composed transcript
(Theorem 3.2, Corollary 5.4). We add only precision: the normalised form
of the floor carries an alphabet-entropy hypothesis that informal
statements frequently elide, and Theorem 3.2 states it.

\textbf{Sequential multi-agent leakage.} The regime this paper addresses
was opened by Asif and Amiri's information-theoretic treatment of agent
pipelines: under their Markov structure, with mutually independent
sensitive variables and marginal per-agent budgets, final-output leakage
is bounded by the geometric sum, and their measurements show mutual
information growing with chain depth {[}asif2026infotheoretic{]}. We
import their Theorem 4.1 with citation, carrying its hypotheses exactly
(Section 4). Our Proposition 4.2 closes the same doubling recursion
inside our model under three named local hypotheses: no carried state,
independent noise, and bounded contamination. The two hypothesis sets
are formally incomparable, neither implying the other: their model
constrains the secrets' joint law and the pipeline's Markov structure,
ours prices the coupling between upstream flow and downstream contexts
directly, and neither result claims the other's hypotheses. AgentLeak
supplies the system-scale measurement: inter-agent channels are
observable in deployed frameworks and carry the dominant share of
leakage, which is what motivates stating our main results against the
transcript adversary (Definition 3.6) {[}elyagoubi2026agentleak{]}.
Patil, Stengel-Eskin and Bansal reach the same qualitative conclusion by
composition analysis {[}patil2025composition{]}. Relative to this line,
our contribution is not a sharper amplification bound; it is the
characterisation of a discipline under which amplification collapses to
additivity, with necessity constructions for both of the discipline's
central hypotheses.

\textbf{Covert channels.} ER-1 (declared channels only) is a
covert-channel exclusion obligation, and its conformance audit is a
channel inventory. The quantitative treatment of unintended channels as
information channels goes back to Millen, who applied Shannon capacity
to covert channels in secure systems {[}millen1987covert{]}. The
necessity construction of Proposition 5.5 can be read in those terms:
state carried across an invocation boundary is an unpriced channel, and
the composition-level bound fails through it at zero declared budget.

\textbf{Language-based information flow.} Static, language-level
enforcement of noninterference-style properties is surveyed by Sabelfeld
and Myers {[}sabelfeld2003language{]}. That programme certifies flows by
analysing program text before execution. The setting here is
complementary: the composed agents are opaque, no per-agent program
analysis is assumed, and the requirements of Section 3.5 are
architectural obligations on the composition, with budgets standing
where typing judgements cannot reach; section 6 records which
requirements are auditable by inspection and which require statistical
estimation. Nothing here replaces language-based enforcement inside an
agent; the results concern what composition preserves when per-agent
internals are inaccessible.

\textbf{Differential privacy composition.} Differential privacy is the
reference discipline for budgets that compose: basic sequential
composition is a theorem (the ε parameters add), advanced composition
degrades sublinearly, and the reason composition holds with no
architectural assumptions is that the guarantee is worst-case over
neighbouring inputs and over adaptively chosen mechanisms
{[}dwork2014algorithmic{]}. That worst-case quantification is the honest
analogue of ER-5's conditioning on the received interface: both price
the adversary's best use of what earlier stages expose. The results here
can be read as identifying what that convenience costs to reproduce for
mutual-information budgets over opaque agents: additivity is not
automatic (Section 4 exhibits exponential amplification, and complete
failure against the transcript adversary at zero budget), and it must be
earned by architecture, namely erasure, conditioning, and provenance
(Theorem 5.1). We make no claim that the budgets of ER-5 imply any
differential-privacy guarantee; the two disciplines quantify over
different adversaries and different neighbouring relations.

\textbf{Time-varying capacity.} Evaluating secrecy against a varying
channel state has precedent in the fading wiretap channel literature
{[}gopala2008secrecy{]}. ER-6's time-indexed budgets and the expiring
deficit of Corollary 5.4 follow that precedent at the system level:
capacity against the strongest available decoder class is a function of
time, and the strict component of any guarantee is dated. The
side-information erosion of Corollary 5.4b is a different time
dependence from the fading channel's: there the channel state varies
while the source is fixed; here every channel and budget is fixed and
what varies is the adversary's background knowledge of the source,
entering through the conditional-entropy term of Fano's inequality
{[}fano1961transmission; cover2006elements{]}.

\subsection{3. Model and preliminaries}\label{model-and-preliminaries}

\subsubsection{3.1 Notation}\label{notation}

Random variables are uppercase, realisations lowercase, alphabets
calligraphic. X is a source (the data subject's protected variable) over
a finite alphabet 𝒳 with H(X) \textgreater{} 0. I(·;·) denotes mutual
information, I(·;·\textbar·) conditional mutual information, H(·)
Shannon entropy in bits. ``U ⊥ V'' abbreviates independence; ``U ⊥ V
\textbar{} W'' conditional independence. For a sequence (A\_1, \ldots,
A\_N) we write A\_\{1:i\} = (A\_1, \ldots, A\_i) and A\_\{\textless i\}
= A\_\{1:i−1\}.

All results in this paper are information-theoretic statements about a
stated model under stated hypotheses. Every hypothesis is listed in
section 3.5 as a numbered engineering requirement or stated locally in
the result that uses it; no result transfers to a system in which the
corresponding hypothesis fails, and section 5 exhibits explicit failure
constructions.

\subsubsection{3.2 The two-agent base
architecture}\label{the-two-agent-base-architecture}

Two agents observe the source through separate channels: a boundary
agent S emitting Y\_S = E\_S(X, N\_S) and a delegation agent M emitting
Y\_M = E\_M(X, N\_M), where N\_S, N\_M are independent local randomness
sources. The architectural intent is that S enforces selective
disclosure under a budget I(X; Y\_S) ≤ C\_S while M executes delegated
capability under a budget I(X; Y\_M) ≤ C\_M, with no unmediated channel
between the two.

The following three results are standard; we state and prove them
in-model because the composition theorems of sections 4 and 5 reduce to
them at N = 2, and because their precise hypotheses are load-bearing.
They are instances of an established family: wiretap equivocation
{[}wyner1975wiretap{]}, secrecy capacity as a capacity difference
{[}leungyancheong1978gaussian{]}, broadcast channels with confidential
messages as the family's general single-eavesdropper form
{[}csiszar1978broadcast{]}, and the source-coding and Fano converses
{[}fano1961transmission; cover2006elements{]}. The failure of the
additive form under collusion needs no citation: it is the elementary
fact that joint information can exceed the sum of marginal informations,
and the interaction identity in Theorem 3.1's proof locates the excess
exactly.

\textbf{Theorem 3.1 (additive bound under conditional independence).} If
I(Y\_S; Y\_M \textbar{} X) = 0, then

\begin{quote}
I(X; Y\_S, Y\_M) ≤ I(X; Y\_S) + I(X; Y\_M),
\end{quote}

with equality if and only if additionally I(Y\_S; Y\_M) = 0.

\emph{Proof.} By the chain rule, I(X; Y\_S, Y\_M) = I(X; Y\_S) + I(X;
Y\_M \textbar{} Y\_S). The interaction identity I(X; Y\_M) − I(X; Y\_M
\textbar{} Y\_S) = I(Y\_S; Y\_M) − I(Y\_S; Y\_M \textbar{} X) (both
sides equal the co-information of the triple) gives, under I(Y\_S; Y\_M
\textbar{} X) = 0,

\begin{quote}
I(X; Y\_M \textbar{} Y\_S) = I(X; Y\_M) − I(Y\_S; Y\_M) ≤ I(X; Y\_M),
\end{quote}

with equality exactly when I(Y\_S; Y\_M) = 0. ∎

The bound is an inequality; conditional independence given X still
permits marginal correlation of Y\_S and Y\_M through X itself, and that
correlation only lowers the joint leakage relative to the sum. The
inequality direction is all that any downstream guarantee uses.

\textbf{Theorem 3.2 (Fano error floor).} Let \textbar 𝒳\textbar{} ≥ 3
and let Y be any observation with I(X; Y) \textless{} ∞. For any
estimator X̂(Y) taking values in 𝒳,

\begin{quote}
P\_e := Pr{[}X̂(Y) ≠ X{]} ≥ (H(X) − I(X; Y) − 1) /
log(\textbar 𝒳\textbar{} − 1).
\end{quote}

If moreover H(X) ≥ log(\textbar 𝒳\textbar{} − 1) (in particular if X is
uniform), then with R := I(X; Y)/H(X),

\begin{quote}
P\_e ≥ 1 − R − 1/H(X).
\end{quote}

\emph{Proof.} Fano's inequality {[}fano1961transmission;
cover2006elements{]} gives H(X \textbar{} Y) ≤ h(P\_e) + P\_e
log(\textbar 𝒳\textbar{} − 1) ≤ 1 + P\_e log(\textbar 𝒳\textbar{} − 1),
where h is binary entropy. Rearranging and substituting H(X \textbar{}
Y) = H(X) − I(X; Y) gives the first display. The second follows since
(H(X) − I − 1)/log(\textbar 𝒳\textbar{} − 1) ≥ (H(X) − I − 1)/H(X) when
H(X) ≥ log(\textbar 𝒳\textbar{} − 1) and H(X) − I − 1 ≥ 0; when H(X) − I
− 1 \textless{} 0 the second display is vacuous and holds trivially. ∎

We record a precision point: the normalised form 1 − R − 1/H(X) requires
the stated alphabet-entropy hypothesis; for strongly non-uniform sources
with H(X) \textless{} log(\textbar 𝒳\textbar{} − 1) only the first
display applies. The log(\textbar 𝒳\textbar{} − 1) denominator
additionally requires the estimator to range over 𝒳 itself; for
estimators taking values in a strict superset of 𝒳 the denominator
weakens to log \textbar 𝒳\textbar.

\textbf{Theorem 3.3 (robustness to approximate separation).} If I(Y\_S;
Y\_M \textbar{} X) ≤ ε, then I(X; Y\_S, Y\_M) ≤ I(X; Y\_S) + I(X; Y\_M)
+ ε.

\emph{Proof.} As in Theorem 3.1, I(X; Y\_M \textbar{} Y\_S) = I(X; Y\_M)
− I(Y\_S; Y\_M) + I(Y\_S; Y\_M \textbar{} X) ≤ I(X; Y\_M) + ε. ∎

Failure of separation is therefore continuous: a violation of size ε
costs at most ε additional bits of joint leakage.

\textbf{Remark 3.4 (the decomposition; what the architecture does and
does not buy).} Write R\_max = (C\_S + C\_M)/H(X). Theorems 3.1 and 3.2
yield: under the separation requirement (ER-1 below) and against a
stated adversary class (ER-6 below), leakage budgets add, and any
estimator faces the error floor P\_e ≥ 1 − R\_max − 1/H(X)
(uniform-source form). These structural requirements do \textbf{not} by
themselves place R\_max below one: two conditionally independent
channels of sufficient combined capacity satisfy every structural
requirement with R\_max ≥ 1, at which point the floor is vacuous. The
strict bound R\_max \textless{} 1 holds exactly when, additionally, the
capacity-deficit condition

\begin{quote}
C\_S + C\_M \textless{} H(X)
\end{quote}

holds. The deficit condition is a measurable, declarable, numerical fact
about a given system and adversary class; it is not a consequence of the
architecture, and we deliberately do not list it among the structural
requirements. It is also a time-indexed quantity: with capacities
evaluated against the strongest decoder class available at time t, the
governing ratio is R(t) = (C\_S(t) + C\_M(t))/H(X) with shelf life t* =
sup\{t : R(t) \textless{} 1\}; R(t) can cross one with every structural
requirement intact, and what expires at t* is the deficit condition, not
the architecture. The supremum is not in general a first crossing: R(t)
may cross one and return, in which case t* marks the last time the
deficit condition holds, not the first time it fails. When C\_S(t) +
C\_M(t) is non-decreasing in t (the regime of ER-6's growing decoder
class, and the hypothesis under which Corollary 5.4 below states the
composed shelf life), R(t) is non-decreasing, the supremum coincides
with the first crossing, and the expiry reading is exact. Treating
channel capacity as a time-indexed quantity has precedent in the study
of fading wiretap channels, where secrecy capacity is defined against a
varying channel state {[}gopala2008secrecy{]}. This decomposition
(structural requirements license additivity and the floor; the declared
deficit licenses the strict bound; the deficit expires) is the paper's
framing discipline, and every bound below is stated within it. A second,
distinct clock, whose drift is real rather than certificational, is
treated at Corollary 5.4b: there the capacities and budgets are held
fixed and what moves is the informed adversary's residual uncertainty
H(X \textbar{} B\_t) about the source. The two clocks must not be
conflated, and every time-indexed reconstruction statement downstream of
this paper should name which one it runs on.

\subsubsection{3.3 N-agent sequential
composition}\label{n-agent-sequential-composition}

We now generalise from two parallel observers to a chain of N
sequentially invoked agents, the regime measured by the recent
multi-agent leakage literature {[}asif2026infotheoretic;
elyagoubi2026agentleak{]}.

\textbf{Definition 3.5 (sequential composition).} A \emph{sequential
composition of N agents} consists of invocations i = 1, \ldots, N where
invocation i:

\begin{enumerate}
\def\labelenumi{\arabic{enumi}.}
\tightlist
\item
  receives an \emph{interface message} Z\_\{i−1\} from its predecessor
  (Z\_0 := ⊥, a constant);
\item
  holds a \emph{sensitive context} S\_i and \emph{carried state}
  ρ\_\{i−1\} (ρ\_0 := ⊥);
\item
  draws \emph{local randomness} N\_i;
\item
  emits a \emph{public output} O\_i, a \emph{forward interface} Z\_i,
  and carried state ρ\_i via a measurable map
\end{enumerate}

\begin{quote}
(O\_i, Z\_i, ρ\_i) = f\_i(Z\_\{i−1\}, ρ\_\{i−1\}, S\_i, N\_i).
\end{quote}

We write V\_i := (O\_i, Z\_i) for the invocation's total emission and T
:= (V\_1, \ldots, V\_N) for the \emph{transcript}.

\textbf{Definition 3.6 (adversary classes).} The \emph{final-output
adversary} 𝔄\_out observes O\_N. The \emph{transcript adversary} 𝔄\_tr
observes T. Clearly any bound against 𝔄\_tr implies the same bound
against 𝔄\_out. Measurement at system scale shows the channels internal
to T are observable in deployed systems and carry the dominant share of
leakage (68.8 per cent of inter-agent channel traces leaking, against
27.2 per cent per-channel output leakage, in a 4,979-trace benchmark
across five production LLMs {[}elyagoubi2026agentleak{]}); we therefore
treat 𝔄\_tr as the realistic passive adversary and state our main
results against it.

\textbf{Definition 3.7 (structural context erasure, as a composition
property).} The composition is \emph{context-erasing at boundary i} if
ρ\_i = ⊥: no state survives the invocation boundary except the declared
forward interface. The composition has \emph{structural context erasure}
if it is context-erasing at every boundary, in which case Definition 3.5
reduces to

\begin{quote}
(O\_i, Z\_i) = f\_i(Z\_\{i−1\}, S\_i, N\_i).
\end{quote}

Consequently (interface sufficiency): for every i, conditional on (Z\_i,
S\_\{i+1\}, N\_\{i+1\}), the emission V\_\{i+1\} is independent of the
entire invocation-i history (Z\_\{i−1\}, S\_i, N\_i). The implication is
one-way: retained state may be causally inert, so the conditional
independence does not entail ρ\_i = ⊥, and erasure is defined by the
structural condition, not by the independence it implies. Erasure is
thus a \emph{reachability} property of the composition operator: no
permitted continuation of the system can reconstruct the erased working
context except through what the declared interface carries, and the
information available downstream about the erased material is bounded
through Z\_i by the data processing inequality.

\textbf{Definition 3.8 (erasure grades).} Erasure at boundary i is
\emph{Grade 1} (hiding) if there exists key material K in the system
(any variable measurable with respect to the system's retained state,
including operator-held state) such that the erased context is
recoverable as a function of (Z\_i, K). It is \emph{Grade 2}
(forgetting) if no such K exists in the closure of the system's
permitted operations. Grade 1 withholds gluing data; Grade 2 destroys
the witness. The closure of the system's permitted operations is left
informal here, and no theorem below depends on the grade criterion
(section 5.3 states this dependence-freedom explicitly). A candidate
formal criterion separating the grades is posed as Open Problem 5.8; the
definition above is operational and suffices for the theorems.

\textbf{Definition 3.9 (background side information; the informed
adversary).} A \emph{background side-information family} is a collection
\{B\_t : t ≥ 0\} of random variables such that (i) for s ≤ t, B\_s is a
measurable function of B\_t (the background only accumulates), and (ii)
for every t the Markov chain B\_t → X → T holds, equivalently I(T; B\_t
\textbar{} X) = 0. The \emph{informed transcript adversary} 𝔄\_tr(t)
observes (T, B\_t). Hypothesis (ii) states that the background carries
information about the source only through the world's dependence on the
source (linkage corpora, auxiliary records, side priors), never through
a tap on the system: for instance a background containing transcript
material, or one correlated with the invocations' internal randomness,
violates it, and no result below applies to such a background (section
5.4); the displayed condition is exact. The index t is the same calendar
time as ER-6's certification index: ``two clocks'' names two quantities
drifting on one calendar, not two time axes. The family is a declared
model of the adversary's accumulation, in the same sense that ER-6's
decoder class is declared; nothing in this paper asserts what any actual
adversary holds.

\subsubsection{3.4 What the per-agent constraint must range
over}\label{what-the-per-agent-constraint-must-range-over}

The composition theorems below turn on two attributes of the
per-invocation leakage constraint, and we fix vocabulary for them now. A
budget is \emph{output-only} if it constrains O\_i alone, and
\emph{interface-inclusive} if it constrains V\_i = (O\_i, Z\_i). A
budget is \emph{marginal} if it is an unconditional mutual information,
and \emph{conditional} if it conditions on the received interface
Z\_\{i−1\}. Section 4 shows that marginal output-only budgets govern
only the final-output adversary and place no nontrivial constraint on
the transcript adversary, even at budget zero. Section 5 shows that
conditional interface-inclusive budgets, combined with structural
context erasure, yield the linear cap against the transcript adversary.

\subsubsection{3.5 Engineering
requirements}\label{engineering-requirements}

Every hypothesis used by any result below is either one of the following
numbered requirements or stated explicitly in the result itself (section
4's marginal-regime results carry local hypotheses in place). Each
requirement is an obligation on an implementation, not a fact about the
world. The necessity constructions of sections 4 and 5 (Propositions
4.3, 5.5 and 5.5b) show what fails when the conditional budget form
(ER-5), erasure (ER-2) or provenance (ER-4) is dropped alone; no
necessity construction is offered for ER-1 or ER-3, whose sufficiency is
claimed only in conjunction with the rest.

\begin{itemize}
\item
  \textbf{ER-1 (declared channels only; non-collusion).} No channel
  connects one invocation to another except the declared components of
  Definition 3.5: the interfaces \{Z\_i\} and the modelled carried state
  \{ρ\_i\}. No undeclared side channel exists. ER-1 does not itself
  constrain the carried state: the destruction of ρ\_i is ER-2's
  obligation alone. Assembly of a downstream context from upstream
  emissions is not a channel in this sense; it is the environment's
  routing, governed by ER-4. At N = 2 (parallel form) this is I(Y\_S;
  Y\_M \textbar{} X) = 0 together with the absence of any third channel
  carrying the inter-agent residue. ER-1 is what the functional form of
  Definition 3.5 encodes; a system with side channels is outside the
  model. That the additive structure fails without ER-1 is elementary,
  not a cited result: joint information exceeds the sum of marginals as
  soon as the observations correlate beyond what the source explains,
  with Theorem 3.3 pricing the excess and Proposition 5.5 exhibiting the
  extremal case at zero marginal budgets; system-scale measurement shows
  deployed multi-agent frameworks routinely violate ER-1 today
  {[}elyagoubi2026agentleak{]}.
\item
  \textbf{ER-2 (structural context erasure).} ρ\_i = ⊥ at every boundary
  (Definition 3.7): the platform destroys all invocation state except
  the declared forward interface. This is a structural property of the
  runtime (state destruction), not an instructed behaviour of the agent.
\item
  \textbf{ER-3 (fresh local randomness).} The local randomness sources
  N\_1, \ldots, N\_N are mutually independent and jointly independent of
  the source and contexts (X, W\_1, \ldots, W\_N), with W\_i as in ER-4.
\item
  \textbf{ER-4 (context provenance).} Each sensitive context is
  assembled from the source, private freshness, and the received
  interface only: S\_i = φ\_i(X, W\_i, Z\_\{i−1\}), where W\_1, \ldots,
  W\_N are mutually independent and jointly independent of X and of
  \{N\_j\}. In particular no downstream context is assembled from
  upstream \emph{emissions} except through the declared (and therefore
  budgeted) interface. Contamination routed through Z\_\{i−1\} is
  permitted, because ER-5 prices it; contamination routed around Z is
  excluded.
\item
  \textbf{ER-5 (conditional, interface-inclusive leakage budgets).} For
  each invocation a declared budget ε\_i ≥ 0 with
\end{itemize}

\begin{quote}
I(O\_i, Z\_i; X \textbar{} Z\_\{i−1\}) ≤ ε\_i.
\end{quote}

The budget covers everything the invocation emits, including the forward
interface, and is conditioned on what the invocation received. It is a
per-invocation, locally defined quantity: it depends only on the joint
law of one invocation's input interface, emission, and the source.
Certification of a declared bound on it is class-and-time-indexed per
ER-6.

\begin{itemize}
\item
  \textbf{ER-6 (stated adversary class, time-indexed).} All budgets and
  capacities are evaluated against a stated decoder class at a stated
  time t; every ε\_i is ε\_i(t), non-decreasing in t under growth of the
  class. No bound in this paper says anything about a decoder class
  stronger than the one declared. (The formalisation of the ordered
  decoder-class family is developed in companion work on the
  time-indexed ratio R(t); here it enters only through the indexing of
  budgets.)

  \emph{Semantics of the time index (certified bounds).} The quantities
  the theorems consume are Shannon functionals of the joint law: I(O\_i,
  Z\_i; X \textbar{} Z\_\{i−1\}) depends on no decoder class and does
  not change as adversaries improve. The declared ε\_i(t) are certified
  upper bounds on these fixed but unknown quantities: a conformance
  audit at time t certifies ε\_i(t) using the estimation and attack
  tooling of the decoder class available at t, and a certificate is only
  as strong as the class that produced it. What is non-decreasing in t
  is therefore the certified value, not the information: growth of the
  class can expose leakage a weaker class missed, forcing the honest
  certificate upward or its revocation, while the underlying mutual
  information is constant. Theorem 5.1 and its corollaries hold with any
  values that genuinely bound the Shannon quantities; every time-indexed
  statement in this paper, in particular the expiry of Corollary 5.4, is
  a statement about what the certification licenses at time t, not a
  statement that information accrues to the transcript. The same reading
  applies to the time-indexed capacities of Remark 3.4. This certified
  clock is not the only time dependence a deployment faces: Corollary
  5.4b treats the informed adversary whose background side information
  grows with calendar time (Definition 3.9), a drift that is real rather
  than certificational and that no conformance audit arrests.
\item
  \textbf{ER-7 (declared deficit for strict reconstruction claims).} Any
  claim that reconstruction is strictly incomplete (error floor bounded
  away from zero) requires, in addition to ER-1 to ER-6, the declared
  numerical fact Σ\_i ε\_i(t) \textless{} H(X) (base case: C\_S + C\_M
  \textless{} H(X)). This is a property of a system-adversary pair at a
  time, never a consequence of ER-1 to ER-5, and it is the component
  that expires as t grows (Remark 3.4).
\end{itemize}

\subsection{4. The sequential bound without
erasure}\label{the-sequential-bound-without-erasure}

This section treats the budget regime the empirical literature measures:
chained agents under marginal output-only budgets. Throughout this
section the composition model is Definition 3.5, the budget is the
marginal output-only constraint I(O\_i; S\_i) ≤ ε\_i on each agent's
leakage about its own context, and the secrets of interest are the
per-agent contexts S\_\{1:N\}. Structural hypotheses are stated per
result and are not uniform across the section: Theorem 4.1 is imported
with its source's Markov structure; Proposition 4.2 assumes no carried
state and independent noise (its hypotheses (i) and (ii) below);
Proposition 4.3's construction satisfies structural context erasure
outright. Carried state and contexts assembled from upstream outputs are
permitted by the model but are excluded wherever a result's stated
hypotheses exclude them.

\textbf{Theorem 4.1 (sequential amplification; imported).} Consider a
sequential pipeline of N agents satisfying the source's Markov structure
(its Eq. (2): S\_\{i−1\} → (O\_\{i−1\}, D\_i, S\_i) → O\_i → O\_\{i+1\}
→ ⋯ → O\_N for each i ∈ \{1, \ldots, N\}, where D\_i denotes the public
or task-relevant inputs to agent i), and assume the sensitive variables
S\_1, \ldots, S\_N are mutually independent. If each agent satisfies the
local leakage constraint I(O\_i; S\_i) ≤ ε\_i, then the final output
satisfies

\begin{quote}
I(O\_N; S\_1, \ldots, S\_N) ≤ Σ\_\{i=1\}\^{}\{N\} 2\^{}\{N−i\} ε\_i,
\end{quote}

and in the uniform case ε\_i = ε,

\begin{quote}
I(O\_N; S\_1, \ldots, S\_N) ≤ (2\^{}N − 1) ε.
\end{quote}

\emph{Provenance.} This is Theorem 4.1 of Asif and Amiri
{[}asif2026infotheoretic{]}, proved there via the chain rule and the
conditional data processing inequality; we import it with citation and
do not re-derive it. The same authors report the empirical counterpart:
average mutual information rising from 0.49 to 1.05 as chain depth grows
from two to five agents (MedQA, LLaMA-7B). Composition analysis by
Patil, Stengel-Eskin and Bansal reaches the same qualitative conclusion
{[}patil2025composition{]}.

\textbf{Proposition 4.2 (in-model reproduction under bounded
contamination).} In the model of Definition 3.5 with O\_i = Z\_i (full
passthrough) and per-agent constraint I(O\_i; S\_i) ≤ ε\_i, suppose
additionally:

\begin{enumerate}
\def\labelenumi{(\roman{enumi})}
\item
  \emph{no carried state}: ρ\_i = ⊥ for every i, equivalently O\_i =
  f\_i(O\_\{i−1\}, S\_i, N\_i); under full passthrough this is
  structural context erasure (Definition 3.7);
\item
  \emph{independent noise}: N\_1, \ldots, N\_N are mutually independent
  and jointly independent of (S\_1, \ldots, S\_N);
\item
  \emph{bounded contamination}: for every i,
\end{enumerate}

\begin{quote}
I(O\_\{i−1\}; S\_i \textbar{} S\_\{1:i−1\}) ≤ I(O\_\{i−1\};
S\_\{1:i−1\}).
\end{quote}

Then L\_i := I(O\_i; S\_\{1:i\}) satisfies the recursion L\_i ≤ 2
L\_\{i−1\} + ε\_i, and hence L\_N ≤ Σ\_\{i=1\}\^{}\{N\} 2\^{}\{N−i\}
ε\_i.

\emph{Proof.} By the chain rule, L\_i = I(O\_i; S\_i) + I(O\_i;
S\_\{1:i−1\} \textbar{} S\_i) ≤ ε\_i + I(O\_i; S\_\{1:i−1\} \textbar{}
S\_i). By (i), O\_i = f\_i(O\_\{i−1\}, S\_i, N\_i), and inductively
O\_\{i−1\} is a measurable function of (S\_\{1:i−1\}, N\_\{1:i−1\});
hence by (ii), N\_i is independent of (S\_\{1:i\}, O\_\{i−1\}) jointly,
and in particular N\_i ⊥ (S\_\{1:i−1\}, O\_\{i−1\}) \textbar{} S\_i.
Given S\_i, the emission O\_i is therefore a measurable function of
(O\_\{i−1\}, N\_i) with N\_i conditionally independent of (S\_\{1:i−1\},
O\_\{i−1\}); hence S\_\{1:i−1\} → O\_\{i−1\} → O\_i is a Markov chain
conditionally on S\_i, and the conditional data processing inequality
gives I(O\_i; S\_\{1:i−1\} \textbar{} S\_i) ≤ I(O\_\{i−1\}; S\_\{1:i−1\}
\textbar{} S\_i). Expanding,

\begin{quote}
I(O\_\{i−1\}; S\_\{1:i−1\} \textbar{} S\_i) = I(O\_\{i−1\};
S\_\{1:i−1\}) + I(O\_\{i−1\}; S\_i \textbar{} S\_\{1:i−1\}) −
I(O\_\{i−1\}; S\_i) ≤ L\_\{i−1\} + I(O\_\{i−1\}; S\_i \textbar{}
S\_\{1:i−1\}) ≤ 2 L\_\{i−1\},
\end{quote}

using bounded contamination in the last step. The recursion L\_i ≤
2L\_\{i−1\} + ε\_i with L\_0 = 0 solves to L\_N ≤ Σ\_i 2\^{}\{N−i\}
ε\_i. ∎

\emph{Scope note (binding).} Bounded contamination is \textbf{our}
sufficient condition for closing the doubling recursion inside this
model; it names the mechanism (the downstream context S\_i correlates
with the upstream transcript at most as strongly as the upstream secrets
do). Whether it coincides with the hypotheses of
{[}asif2026infotheoretic{]} is a citation-resolution item; Theorem 4.1
stands on its own citation and Proposition 4.2 on its own proof, and
neither claims the other's hypotheses.

The doubling term is genuinely a contamination term: when the downstream
context is independent of the upstream flow given earlier secrets
(I(O\_\{i−1\}; S\_i \textbar{} S\_\{1:i−1\}) = 0, the exogenous case),
the same computation yields the additive recursion L\_i ≤ L\_\{i−1\} +
ε\_i under the same no-carried-state hypothesis, \emph{for the
final-output adversary}. Hypothesis (i) is load-bearing, not decorative:
with carried state permitted, zero marginal budgets bound nothing even
at the final output. (Take S\_1, S\_2 independent uniform bits;
invocation 1 emits fresh junk O\_1 = U with U uniform and independent of
the secrets, retaining ρ\_1 = S\_1; invocation 2 emits O\_2 = ρ\_1.
Every marginal budget holds at ε = 0 and bounded contamination holds
with both sides zero, yet I(O\_2; S\_1, S\_2) = H(S\_1).) The regimes
are therefore three: with carried state, no function of the marginal
budgets bounds leakage (here at the final output, and at transcript
scale in Proposition 5.5); without carried state, under marginal budgets
only, amplification up to the geometric sum (Proposition 4.2); under
erasure with the conditional budgets of ER-5, the linear cap (Theorem
5.1). The next proposition shows why the final-output adversary is the
wrong one to design against even inside the middle regime.

\textbf{Proposition 4.3 (marginal budgets do not bind the transcript
adversary).} There is a two-agent sequential composition satisfying
structural context erasure (ER-2), fresh randomness (ER-3), exogenous
contexts, and \emph{zero} marginal budgets, in which the transcript
adversary recovers a context completely:

\begin{quote}
I(O\_1; S\_1) = I(O\_2; S\_2) = 0, and moreover I(V\_1; S\_1) = I(V\_2;
S\_2) = 0, yet I(O\_1, O\_2; S\_1, S\_2) = H(S\_2).
\end{quote}

\emph{Construction.} Let S\_1, S\_2 be independent uniform bits and N\_1
a uniform bit independent of both. Agent 1 emits O\_1 = Z\_1 = S\_1 ⊕
N\_1. Agent 2 emits O\_2 = Z\_1 ⊕ S\_2 (and Z\_2 = ⊥). Marginally, O\_1
is uniform and independent of S\_1 (one-time pad), and O\_2 = S\_1 ⊕
N\_1 ⊕ S\_2 is uniform and independent of S\_2; so every marginal budget
is met at ε = 0, including the interface-inclusive marginal I(V\_i;
S\_i) = 0. But O\_1 ⊕ O\_2 = S\_2: the transcript adversary reads S\_2
exactly. Consistently with Theorem 4.1, the final output alone reveals
nothing (I(O\_2; S\_1, S\_2) = 0). ∎

Three consequences. First, a bound stated against the final output, such
as Theorem 4.1, does not survive the passage to the transcript
adversary, and the transcript adversary is the one system-scale
measurement says exists {[}elyagoubi2026agentleak{]}. Second, the
failure is not repaired by including the interface in a \emph{marginal}
budget; the construction defeats that too. The load-bearing attribute is
\emph{conditioning}: the conditional budget of ER-5 evaluates I(V\_2;
S\_2 \textbar{} Z\_1) = H(S\_2) and prices the correlation the marginal
misses. Third, the construction satisfies erasure: structural context
erasure alone, without the conditional budget discipline, caps nothing.
Erasure and budget form are separately necessary; section 5 shows they
are jointly sufficient. These three facts make Proposition 4.3 the
paper's unconditional separation: zero declared budget on one side
against total recovery of a context on the other, with no attainability
hypothesis on either side. The exponential-to-linear comparison of
Remark 5.7 is, by contrast, a comparison of upper-bound guarantees under
stated budget semantics (see the second fence there).

\subsection{5. The linear cap under structural context
erasure}\label{the-linear-cap-under-structural-context-erasure}

\subsubsection{5.1 The cap}\label{the-cap}

\textbf{Theorem 5.1 (linear cap).} Consider a sequential composition of
N agents satisfying ER-1 (declared channels only), ER-2 (structural
context erasure), ER-3 (fresh randomness), ER-4 (context provenance),
and ER-5 (conditional interface-inclusive budgets ε\_1, \ldots, ε\_N).
Then against the transcript adversary,

\begin{quote}
I(T; X) ≤ Σ\_\{i=1\}\^{}\{N\} ε\_i,
\end{quote}

and in the uniform case I(T; X) ≤ Nε.

\emph{Proof.} Write V\_i = (O\_i, Z\_i) and T = (V\_1, \ldots, V\_N). By
the chain rule,

\begin{quote}
I(T; X) = Σ\_\{i=1\}\^{}\{N\} I(V\_i; X \textbar{} V\_\{1:i−1\}).
\end{quote}

Fix i and set A := V\_\{1:i−1\}; note Z\_\{i−1\} is a coordinate of A
(for i = 1, both are trivial).

\emph{Step 1 (transcript measurability).} By ER-1, ER-2 and ER-4, each
V\_j with j \textless{} i is a measurable function of (X, W\_\{1:j\},
N\_\{1:j\}): inductively, V\_j = f\_j(Z\_\{j−1\}, φ\_j(X, W\_j,
Z\_\{j−1\}), N\_j) and Z\_\{j−1\} is a coordinate of V\_\{j−1\}. Hence A
is a measurable function of (X, W\_\{\textless i\}, N\_\{\textless i\}).

\emph{Step 2 (fresh inputs are conditionally independent of the past).}
By ER-3 and ER-4, (W\_i, N\_i) ⊥ (X, W\_\{\textless i\},
N\_\{\textless i\}), hence by Step 1, (W\_i, N\_i) ⊥ (A, X) jointly.
Since (X, Z\_\{i−1\}) is a measurable function of (A, X), it follows
that (W\_i, N\_i) ⊥ A \textbar{} (X, Z\_\{i−1\}). Conditional on (X,
Z\_\{i−1\}), the emission V\_i = f\_i(Z\_\{i−1\}, φ\_i(X, W\_i,
Z\_\{i−1\}), N\_i) is a measurable function of (W\_i, N\_i) and the
conditioned values, so by the data processing inequality

\begin{quote}
I(V\_i; A \textbar{} X, Z\_\{i−1\}) ≤ I(W\_i, N\_i; A \textbar{} X,
Z\_\{i−1\}) = 0.
\end{quote}

\emph{Step 3 (exchange).} Expanding I(V\_i; (X, A) \textbar{}
Z\_\{i−1\}) by the chain rule in both orders,

\begin{quote}
I(V\_i; X \textbar{} Z\_\{i−1\}) + I(V\_i; A \textbar{} X, Z\_\{i−1\}) =
I(V\_i; A \textbar{} Z\_\{i−1\}) + I(V\_i; X \textbar{} A, Z\_\{i−1\}).
\end{quote}

Since Z\_\{i−1\} is a function of A, I(V\_i; X \textbar{} A, Z\_\{i−1\})
= I(V\_i; X \textbar{} A). With Step 2 this gives

\begin{quote}
I(V\_i; X \textbar{} A) = I(V\_i; X \textbar{} Z\_\{i−1\}) − I(V\_i; A
\textbar{} Z\_\{i−1\}) ≤ I(V\_i; X \textbar{} Z\_\{i−1\}) ≤ ε\_i,
\end{quote}

using ER-5. Summing over i completes the proof. ∎

The proof isolates where each requirement acts. ER-2 (with ER-1) is what
makes the received interface Z\_\{i−1\} a sufficient statistic of the
past for invocation i, so that the conditional budget prices everything
the invocation can correlate with; Step 2 is exactly the point that
fails when carried state survives (Proposition 5.5). ER-4 is what
prevents the environment from smuggling upstream emissions into
downstream contexts around the priced interface. ER-5's conditioning is
what Proposition 4.3 showed to be irreplaceable.

\textbf{Corollary 5.2 (per-agent secrets variant).} Suppose instead of
ER-4 the contexts S\_1, \ldots, S\_N are mutually independent and
jointly independent of \{N\_j\} (exogenous per-agent secrets), and the
budgets take the form I(V\_i; S\_i \textbar{} Z\_\{i−1\}) ≤ ε\_i. Then

\begin{quote}
I(T; S\_1, \ldots, S\_N) ≤ Σ\_\{i=1\}\^{}\{N\} ε\_i.
\end{quote}

\emph{Proof.} Apply Theorem 5.1 with X := (S\_1, \ldots, S\_N), φ\_i the
i-th projection and W\_i trivial; its hypotheses hold. It remains to
reduce the theorem's budget to the stated one:

\begin{quote}
I(V\_i; S\_\{1:N\} \textbar{} Z\_\{i−1\}) = I(V\_i; S\_i \textbar{}
Z\_\{i−1\}) + I(V\_i; S\_\{≠i\} \textbar{} Z\_\{i−1\}, S\_i),
\end{quote}

and the second term vanishes: conditional on (Z\_\{i−1\}, S\_i), V\_i is
a function of N\_i, while N\_i ⊥ (S\_\{1:N\}, N\_\{\textless i\})
jointly and (Z\_\{i−1\}, S\_i, S\_\{≠i\}) are functions of that tuple,
so N\_i ⊥ S\_\{≠i\} \textbar{} (Z\_\{i−1\}, S\_i). ∎

\textbf{Example 5.3 (the two-agent base as an instance).} The parallel
architecture of section 3.2 is the N = 2 composition with Z\_1 = ⊥ and
S\_i = X: the budgets of ER-5 reduce to I(Y\_S; X) ≤ C\_S and I(Y\_M; X)
≤ C\_M, and Theorem 5.1 returns I(Y\_S, Y\_M; X) ≤ C\_S + C\_M, the
additive bound of Theorem 3.1 in budget form. The composition framework
strictly generalises the base architecture.

\textbf{Corollary 5.4 (composed error floor, and the deficit
condition).} Under the hypotheses of Theorem 5.1 with
\textbar 𝒳\textbar{} ≥ 3, any estimator X̂(T) of the source from the full
transcript satisfies

\begin{quote}
P\_e ≥ (H(X) − Σ\_i ε\_i − 1) / log(\textbar 𝒳\textbar{} − 1),
\end{quote}

and, for H(X) ≥ log(\textbar 𝒳\textbar{} − 1), P\_e ≥ 1 − (Σ\_i ε\_i +
1)/H(X). Per ER-7, the declared deficit Σ\_i ε\_i \textless{} H(X)
licenses strict incompleteness: retaining the binary-entropy term in
Fano's inequality, P\_e = 0 would force H(X \textbar{} T) = 0 and hence
I(T; X) = H(X), contradicting Theorem 5.1 under the deficit; so Σ\_i
ε\_i \textless{} H(X) already implies P\_e \textgreater{} 0. The
displayed floor, which relaxes the binary-entropy term to one bit, is
positive exactly under the stronger condition Σ\_i ε\_i \textless{} H(X)
− 1; between the two thresholds the deficit still licenses P\_e
\textgreater{} 0 while the displayed linear bound is vacuous. The
deficit is a numerical fact about the system-adversary pair at the
evaluation time t of ER-6, not a consequence of ER-1 to ER-5. With
certified budgets ε\_i(t) non-decreasing in t, the composed shelf life
is t*\_lin = sup\{t : Σ\_i ε\_i(t) \textless{} H(X)\}, and what expires
at t*\_lin is the deficit, not the composition discipline: every
structural requirement can hold while the floor decays to nothing. Per
ER-6's certified-bound semantics, the expiry is a claim about the
certification: at t*\_lin the certificates no longer witness a deficit,
and the floor that decays is the certified floor; nothing here asserts
that the transcript's information about the source has grown.

\emph{Proof.} Theorem 5.1 and Theorem 3.2 with Y := T. ∎

\textbf{Corollary 5.4b (informed-adversary floor; the erosion clock).}
Under the hypotheses of Theorem 5.1 with \textbar 𝒳\textbar{} ≥ 3, let
\{B\_t\} be a background side-information family (Definition 3.9). Then:

\begin{enumerate}
\def\labelenumi{(\roman{enumi})}
\tightlist
\item
  \emph{(the cap survives conditioning)} for every t,
\end{enumerate}

\begin{quote}
I(T; X \textbar{} B\_t) = I(T; X) − I(T; B\_t) ≤ I(T; X) ≤ Σ\_i ε\_i;
\end{quote}

\begin{enumerate}
\def\labelenumi{(\roman{enumi})}
\setcounter{enumi}{1}
\tightlist
\item
  \emph{(the floor erodes through the residual entropy)} any estimator
  X̂(T, B\_t) of the informed transcript adversary 𝔄\_tr(t), taking
  values in 𝒳, satisfies
\end{enumerate}

\begin{quote}
P\_e(t) ≥ (H(X \textbar{} B\_t) − Σ\_i ε\_i − 1) /
log(\textbar 𝒳\textbar{} − 1),
\end{quote}

and, when H(X \textbar{} B\_t) ≥ log(\textbar 𝒳\textbar{} − 1), P\_e(t)
≥ 1 − R\_inf(t) − 1/H(X \textbar{} B\_t), where R\_inf(t) := Σ\_i ε\_i /
H(X \textbar{} B\_t);

\begin{enumerate}
\def\labelenumi{(\roman{enumi})}
\setcounter{enumi}{2}
\tightlist
\item
  \emph{(monotonicity and expiry)} H(X \textbar{} B\_t) is
  non-increasing in t. Strict incompleteness of reconstruction by the
  informed adversary is licensed by the \emph{informed deficit
  condition} Σ\_i ε\_i \textless{} H(X \textbar{} B\_t): retaining the
  binary-entropy term in Fano's inequality exactly as at Corollary 5.4,
  P\_e(t) = 0 would force I(T, B\_t; X) = H(X), while I(T, B\_t; X) =
  I(B\_t; X) + I(T; X \textbar{} B\_t) \textless{} (H(X) − H(X
  \textbar{} B\_t)) + H(X \textbar{} B\_t) = H(X) under the informed
  deficit. With certified budgets ε\_i(t) non-decreasing (ER-6) and H(X
  \textbar{} B\_t) non-increasing, the informed shelf life t*\_inf =
  sup\{t : Σ\_i ε\_i(t) \textless{} H(X \textbar{} B\_t)\} coincides
  with the first crossing.
\end{enumerate}

\emph{Proof.} (i) Expanding I(T; X, B\_t) by the chain rule in both
orders, I(T; X) + I(T; B\_t \textbar{} X) = I(T; B\_t) + I(T; X
\textbar{} B\_t); the second term on the left vanishes by Definition
3.9(ii), and Theorem 5.1 bounds I(T; X). (ii) Theorem 3.2 with Y := (T,
B\_t), using H(X) − I(X; T, B\_t) = H(X \textbar{} T, B\_t) = H(X
\textbar{} B\_t) − I(T; X \textbar{} B\_t) ≥ H(X \textbar{} B\_t) − Σ\_i
ε\_i by (i); the normalised form follows as at Theorem 3.2 with H(X
\textbar{} B\_t) in place of H(X). (iii) For s ≤ t, B\_s is a measurable
function of B\_t, so H(X \textbar{} B\_t) = H(X \textbar{} B\_t, B\_s) ≤
H(X \textbar{} B\_s); the exact-Fano computation is displayed in the
statement; a non-decreasing left-hand side against a non-increasing
right-hand side makes the deficit set an interval whose supremum is the
first exit. ∎

The proof uses only Theorem 5.1's conclusion, so the corollary is
portable to any composition carrying a certified transcript cap Σ\_i
ε\_i, however obtained. Three readings fence the corollary. First, what
erodes and what does not. The transcript's leakage about the source does
not grow; by (i) its incremental value to the informed adversary, I(T; X
\textbar{} B\_t) = I(T; X) − I(T; B\_t), is non-increasing in t, since
I(T; B\_t) is non-decreasing by the data processing inequality along the
accumulation. What erodes is the height of the floor, because the
informed adversary's uncertainty about the source starts from H(X
\textbar{} B\_t), not H(X). The protection does not leak more; it
matters less. Second, the clock separation. ER-6's certified drift is
epistemic: it records what an audit can license, and better estimation
tooling can arrest or even reverse an over-conservative certificate. The
drift here is real in the model's sense: what moves is the modelled
quantity itself, not a certificate of a fixed quantity; the adversary's
residual uncertainty does not rise, any fall in it is genuine, and no
audit arrests it. A deployed guarantee therefore carries two dated
components, a certified deficit against a declared decoder class
(Corollary 5.4) and an informed deficit against a declared background
family (this corollary), and the two expire independently. Third,
vacuity at the far end. The normalised form's hypothesis H(X \textbar{}
B\_t) ≥ log(\textbar 𝒳\textbar{} − 1) itself expires as the background
accumulates, leaving only the unnormalised display; and past t*\_inf the
floor is vacuous for every architecture satisfying every requirement,
because the corollary bounds what the transcript adds to an informed
adversary and cannot prevent the world from learning the source by other
means.

\subsubsection{5.2 Necessity and
tightness}\label{necessity-and-tightness}

\textbf{Proposition 5.5 (necessity of erasure).} Drop ER-2 only. Then no
bound of the form I(T; X) ≤ g(Σ\_i ε\_i) with g(0) = 0 holds: there is a
two-invocation composition with carried state, satisfying ER-1, ER-3,
ER-4 and ER-5 with ε\_1 = ε\_2 = 0, in which I(T; X) = H(X).

\emph{Construction.} Let X be uniform on \{0,1\}\^{}m and let the agent
draw a persistent uniform pad ρ ∈ \{0,1\}\^{}m at first invocation (ρ ⊥
X), retaining it as carried state. Invocation 1 emits V\_1 = O\_1 = X ⊕
ρ with Z\_1 = ⊥; invocation 2 emits V\_2 = O\_2 = ρ. Both conditional
budgets are zero: I(X ⊕ ρ; X \textbar{} ⊥) = 0 and I(ρ; X \textbar{} ⊥)
= 0. Yet O\_1 ⊕ O\_2 = X, so I(T; X) = m = H(X), with m arbitrary. The
proof of Theorem 5.1 fails exactly at Step 2: with ρ surviving the
boundary, V\_2 is not conditionally independent of V\_1 given (X, Z\_1),
and the correlation carried by ρ is priced by no budget. ∎

Each emission is individually a one-time-pad ciphertext or a pad,
individually independent of the source; jointly they determine it. The
construction is the composition-level form of the classical observation
that per-component independence does not bound joint information, and it
identifies undestroyed cross-invocation state as an unpriced channel.
Together with Proposition 4.3 (which showed the conditional budget
necessary while erasure held), it shows ER-2 and ER-5 are individually
necessary; the next proposition adds the third construction, for ER-4.

\textbf{Proposition 5.5b (necessity of provenance).} Drop ER-4 only.
Then no bound of the form I(T; X) ≤ g(Σ\_i ε\_i) with g(0) = 0 holds:
there is a two-invocation composition satisfying ER-1, ER-2, ER-3 and
ER-5 with ε\_1 = ε\_2 = 0, in which a downstream context is assembled
from an upstream emission around the declared interface and I(T; X) =
H(X).

\emph{Construction.} Let X be uniform on \{0,1\}\^{}m. Invocation 1
holds a trivial context (φ\_1 constant), draws N\_1 uniform on
\{0,1\}\^{}m independent of X, and emits V\_1 = O\_1 = N\_1 with Z\_1 =
⊥ and ρ\_1 = ⊥. The environment assembles the downstream context from
the public upstream emission: S\_2 = X ⊕ O\_1, which violates ER-4
(assembly from an emission other than through Z\_1). Invocation 2 emits
V\_2 = O\_2 = S\_2 with Z\_2 = ⊥ and ρ\_2 = ⊥. Both conditional
interface-inclusive budgets are zero: I(V\_1; X \textbar{} Z\_0) =
I(N\_1; X) = 0 and I(V\_2; X \textbar{} Z\_1) = I(X ⊕ N\_1; X) = 0. Yet
O\_1 ⊕ O\_2 = X, so I(T; X) = m = H(X). Erasure holds at every boundary
and no undeclared channel connects the invocations; the leak travels
entirely through the environment's context assembly, the route ER-4
exists to exclude. The proof of Theorem 5.1 fails exactly at Step 2:
V\_2 is not of the form f\_2(Z\_1, φ\_2(X, W\_2, Z\_1), N\_2), and its
dependence on the upstream randomness N\_1 is priced by no conditional
budget, since Z\_1 = ⊥ carries nothing to condition on. ∎

Propositions 4.3, 5.5 and 5.5b are the paper's necessity constructions:
the conditional form of the budget (ER-5), structural context erasure
(ER-2), and context provenance (ER-4) are individually indispensable,
each defeated at declared budget zero when dropped alone. No necessity
construction is offered for ER-1 (declared channels) or ER-3 (fresh
randomness); the results claim their sufficiency in conjunction with the
rest, not their individual indispensability. By Theorem 5.1 the five
together are sufficient.

\textbf{Proposition 5.6 (tightness).} The cap of Theorem 5.1 is
achieved: for X = (X\_1, \ldots, X\_N) with independent uniform bits,
the composition O\_i = X\_i, Z\_i = ⊥ satisfies all hypotheses with ε\_i
= 1 and I(T; X) = N = Σ\_i ε\_i. Hence Nε cannot be improved without
further hypotheses.

\textbf{Remark 5.7 (the exponential-to-linear separation, and what
remains empirical).} Against the final-output adversary with marginal
budgets and contaminated contexts, the available guarantee is the
geometric sum Σ 2\^{}\{N−i\} ε\_i, i.e.~(2\^{}N − 1)ε uniformly (Theorem
4.1); under the composition discipline ER-1 to ER-5 the guarantee
against the transcript adversary is Nε (Theorem 5.1). At N = 2 the
uniform comparison is 3ε against 2ε; at N = 5, 31ε against 5ε; the ratio
(2\^{}N − 1)/N grows without bound. Three fences on interpreting the
separation. First, it is a separation between \emph{disciplines}, not
mechanisms: Proposition 4.3 shows erasure without conditional
interface-inclusive budgets caps nothing, Proposition 5.5 shows budgets
without erasure cap nothing, and Proposition 5.5b shows budgets with
erasure but without provenance cap nothing; the linear regime requires
the conjunction. Second, it is a comparison of \emph{guarantees}, not of
attainable leakages, and its two ε parameters are not the same object.
Both displays are upper bounds: no construction here or in the cited
literature attains leakage of geometric order under marginal budgets,
and the source's own depth measurements (Theorem 4.1's provenance note)
do not exhibit geometric growth. Moreover the conditional
interface-inclusive budget at level ε is the strictly more demanding
constraint on an implementation than a marginal budget at the same
numeral: Proposition 4.3's composition meets every marginal budget at
zero while its conditional budget is I(V\_2; S\_2 \textbar{} Z\_1) =
H(S\_2). The comparison is therefore between the best guarantees
available under each stated budget semantics, not a like-for-like ratio
of leakages; the unconditional separation this paper proves is
Proposition 4.3, where zero declared budget coexists with total
transcript recovery of a context. Third, the mathematical content is now
exhausted by the conditional theorems: what remains open is empirical,
namely whether any given deployed system satisfies ER-1 to ER-5. That is
a conformance property of an implementation, testable per instance by
channel inventory (ER-1), runtime state audit (ER-2, ER-3), provenance
analysis (ER-4) and leakage measurement against the declared class
(ER-5, ER-6); it is an engineering hypothesis about a system, never a
theorem about all systems, and this paper proves nothing about
implementations that have not been so audited.

\subsubsection{5.3 The grade criterion, stated as an open
problem}\label{the-grade-criterion-stated-as-an-open-problem}

The theorems above use the operational grade distinction of Definition
3.8 and never the deeper criterion. The deeper criterion is open, and we
state it as an open problem with its proof obligations, because a
positive resolution upgrades the \emph{time} status of Grade-2 erasure:
a Grade-2-erased context leaves nothing archived for a stronger decoder
class to revisit, making erasure the only mechanism in the model whose
guarantee would carry no certification index (contrast Corollary 5.4,
where the certified deficit expires, and Corollary 5.4b, whose erosion
clock is unaffected either way; see the scope note below).

\textbf{Open Problem 5.8 (obstruction-theoretic grade criterion).} There
exists a formalisation of the N-invocation system in which the retained
local views (the variables surviving each boundary) are sections of a
presheaf of partial reconstructions over a suitable cover of the joint
observation space, such that: (a) erasure at a boundary is Grade 2
(Definition 3.8) if and only if the obstruction class to gluing the
local sections into a global section witnessing the erased context is
non-vanishing; and (b) erasure is Grade 1 if and only if the class
vanishes and only the gluing data (key material) is withheld.

\emph{Proof obligations.} (i) Construct the site: a cover and topology
on which the candidate Čech class lives, for at least the two-agent
instantiation. (ii) Prove the correspondence between
recoverability-with-keys and vanishing of the class. (iii) Exhibit the
class as a computable invariant in one non-trivial instance.

\emph{The missing step, named.} Obligation (i) is unstarted: no
candidate construction of the cover and topology has been written down;
the sheaf-theoretic vocabulary is at present imported, not constructed.
Nothing in sections 3 to 5 depends on this problem.

\emph{Scope of the upgrade.} The time-status upgrade a positive
resolution would license concerns the erasure guarantee itself: nothing
archived remains for a stronger decoder class to revisit, so that
guarantee carries no certification index. It does not arrest the erosion
clock of Corollary 5.4b, which runs on background the adversary
accumulates outside the system and is indifferent to every property of
the composition, erasure included.

\emph{Falsifier scope.} The statement is existentially quantified over
formalisations and is therefore a research programme rather than a
falsifiable conjecture: a failed candidate site refutes that candidate,
not the existential. What is falsifiable is any fixed candidate
criterion: once a site is fixed, any composition of retained views,
under any future capability, that recovers Grade-2-erased material would
exhibit a vanishing effective obstruction where the candidate criterion
demanded a non-vanishing one, refuting that candidate. Obligation (i) is
what converts the problem into a falsifiable conjecture, and it is
unstarted; the label ``open problem'' rather than ``conjecture'' records
exactly this.

\subsubsection{5.4 Limits of the results}\label{limits-of-the-results}

\emph{(Seeds section 6; A3 content, A2 may reframe prose.)}

\begin{itemize}
\tightlist
\item
  Every result is conditional on ER-1 to ER-5, which are obligations,
  not observations. System-scale measurement of deployed multi-agent
  frameworks shows ER-1 is routinely violated today
  {[}elyagoubi2026agentleak{]}; a system that has not been audited
  against the requirements inherits nothing from Theorem 5.1.
\item
  Nothing here bounds active adversaries, adversaries that corrupt an
  agent, or side channels outside the declared interfaces (timing,
  resource usage); the model is passive-observational.
\item
  The cap Nε is a leakage bound, not a reconstruction impossibility.
  Strict incompleteness of reconstruction requires the declared deficit
  of ER-7, which is time-indexed and expires (Corollary 5.4). No
  sentence in this paper asserts a reconstruction ceiling without its
  preconditions and its time index, and none should survive editing that
  does.
\item
  Theorem 4.1 is imported, not re-derived; Proposition 4.2 closes the
  same recursion in-model under local hypotheses (no carried state,
  independent noise, bounded contamination) whose correspondence to the
  source's hypotheses is a citation-resolution obligation.
\item
  The exponential-to-linear comparison is a comparison of upper-bound
  guarantees under different budget semantics (Remark 5.7, second
  fence); no attainability construction for the geometric regime is
  exhibited or claimed. The unconditional separation is Proposition 4.3.
\item
  The budgets ε\_i(t) are certified upper bounds on fixed Shannon
  quantities (ER-6); every expiry statement is a claim about the
  certification, not about information accruing to the transcript
  (Corollary 5.4).
\item
  The informed-adversary floor (Corollary 5.4b) is conditional on the
  Markov hypothesis of Definition 3.9. A background corpus containing
  transcript material violates it; conditioning on such a background can
  increase leakage beyond the cap, and the corollary is then
  inapplicable. The background family is a declared model, in the same
  sense that ER-6's decoder class is declared, and the empirical
  schedule of H(X \textbar{} B\_t) is not estimated here. Past t*\_inf
  the floor is vacuous for every architecture: the corollary bounds what
  the transcript adds to an informed adversary; it cannot prevent the
  world from learning the source by other means.
\item
  The Fano floor's normalised form carries an alphabet-entropy
  hypothesis (Theorem 3.2) frequently elided in informal statements.
\item
  Per the programme's reporting discipline, a measured result against
  these predictions (for instance, a conformant system exhibiting
  super-linear transcript leakage) would be reported with the same
  prominence as a confirmation.
\end{itemize}

\subsection{6. Threats to validity}\label{threats-to-validity}

\textbf{Conditionality and conformance.} Every result in this paper is
conditional on ER-1 to ER-5, which are obligations on implementations,
not observations about the world. System-scale measurement indicates
that deployed multi-agent frameworks routinely violate ER-1 today
{[}elyagoubi2026agentleak{]}. Conformance is testable per instance:
channel inventory for ER-1, runtime state audit for ER-2 and ER-3,
provenance analysis for ER-4, and leakage measurement against the
declared decoder class for ER-5 and ER-6. The requirements are not
equally auditable, and the ER-5 leg is the hard one. Conformance to ER-5
requires estimating a conditional mutual information between an
invocation's full emission and the source given the received interface,
a notoriously hard statistical problem for high-dimensional text
channels, and this paper specifies no estimator. Such estimation is
feasible only against restricted decoder classes, so an ER-5 conformance
verdict is always relative to the class whose estimation and attack
tooling produced it. This is one reason ER-6 indexes budgets by decoder
class and time: under its certified-bound semantics, a certificate for
ε\_i(t) is only as strong as the class that produced it, and growth of
the class can force the honest certificate upward. Conformance is an
engineering hypothesis about a system, never a theorem about all
systems; a system that has not been audited against the requirements
inherits nothing from Theorem 5.1, and this paper asserts conformance
for no deployed system.

\textbf{Adversary model.} The model is passive-observational. Nothing
here bounds active adversaries, adversaries that corrupt or replace an
agent, or side channels outside the declared interfaces, including
timing and resource usage. Extending the ER-1 inventory to physical side
channels is an audit problem this paper does not solve.

\textbf{Leakage versus reconstruction, and expiry.} The cap Σ\_i ε\_i is
a leakage bound, not a reconstruction impossibility. The strict claim
that reconstruction is incomplete additionally requires the declared
capacity-deficit condition of ER-7, evaluated against the stated
adversary class at a stated time, and that condition expires in the
certified sense of ER-6: the certified budgets ε\_i(t) are
non-decreasing as the decoder class grows, and the floor of Corollary
5.4 decays to vacuity at the shelf life t*\_lin while every structural
requirement remains intact. What expires is what the certification
licenses at time t; the transcript's mutual information about the source
is a fixed functional of the joint law, and nothing here asserts that it
grows. No sentence in this paper asserts strict incompleteness of
reconstruction without the deficit condition and its time index, and
none should survive editing that does.

\textbf{Two clocks, and the corpus model.} The expiry of the previous
paragraph is certificational. Corollary 5.4b adds a second, real expiry:
against an informed adversary whose background side information
accumulates, the error floor erodes through H(X \textbar{} B\_t) with
every certificate intact and no growth in what the transcript reveals.
That result is conditional on the declared background family and its
Markov hypothesis (Definition 3.9): whether any actual adversary's
corpus satisfies the hypothesis, and how fast H(X \textbar{} B\_t)
falls, are empirical questions this paper does not answer, symmetric to
the conformance questions above. A deployed guarantee should declare its
corpus model alongside its decoder class, and its two deficits expire
independently.

\textbf{Import fidelity.} Theorem 4.1 is imported on its citation and
not re-derived; its hypotheses (the source's Markov structure and the
mutual independence of the sensitive variables) are carried inside the
imported statement. Proposition 4.2 is proved here under our own local
hypotheses (no carried state, independent noise, and bounded
contamination), and its named condition and the source's hypothesis set
are formally incomparable: the source's hypotheses do not imply bounded
contamination, and bounded contamination does not imply them. No result
in Section 5 depends on that correspondence.

\textbf{Statement precision.} The normalised form of the Fano floor
requires the alphabet-entropy hypothesis of Theorem 3.2; for strongly
non-uniform sources only the unnormalised display applies. The additive
two-agent bound is an inequality with a stated equality condition, not
an identity (Theorem 3.1).

\textbf{Reporting commitment.} A measured result against these
predictions, for instance a system audited as conformant to ER-1 to ER-5
exhibiting super-linear transcript leakage, would falsify the
conformance audit, the model's fit to that system, or the applicability
of the theorem, and it would be reported with the same prominence as a
confirmation.

\subsection{7. Conclusion}\label{conclusion}

We have given a composition model for sequentially invoked agents in
which the two regimes of the recent multi-agent leakage literature are
two sides of one architectural line. Without erasure and conditional
interface-inclusive budgets, per-agent certification does not compose:
the guarantee available on leakage about the joint secrets at the final
output is only the geometric sum (Theorem 4.1, imported; an upper bound,
with no attainability construction known), and against the transcript
adversary marginal certification is vacuous even at budget zero
(Proposition 4.3, the paper's unconditional separation). Under the five
requirements ER-1 to ER-5, transcript leakage is capped by the sum of
the declared budgets (Theorem 5.1); three of the five hypotheses are
individually necessary (Propositions 4.3, 5.5 and 5.5b), and the cap is
tight (Proposition 5.6). The exponential-to-linear comparison is a
comparison of upper-bound guarantees under stated budget semantics, not
of attainable leakages (Remark 5.7).

The framing discipline carries as much weight as the mathematics. What
the architecture guarantees is additive accounting and an error floor
whose height depends on declared numbers. The strict claim that
reconstruction is incomplete is licensed only by a declared
capacity-deficit at a stated time against a stated decoder class, and
that component expires. Corollary 5.4b adds the second clock: even with
every certificate intact, the floor an informed adversary faces is
non-increasing, eroding as background side information accumulates, on a
drift no audit arrests. Guarantees for deployed systems should be stated
in this decomposed form: structural properties audited, numerical
conditions declared and dated, the decoder class and the background
corpus model both named.

Three items remain open. First, conformance: whether any given deployed
system satisfies ER-1 to ER-5 is an empirical, per-instance question,
and benchmark measurement of candidate conformant systems is future
work. Second, the grade criterion: Open Problem 5.8 states a candidate
formal separation between erasure that hides and erasure that forgets,
with its proof obligations and its unstarted first step named; a
positive resolution would upgrade the time status of the strongest
erasure grade, since nothing archived would remain for a stronger
decoder class to revisit. Third, the adversary model: active adversaries
and physical side channels are outside the model, and the requirements'
audit surface for them is undefined. The results are offered in the
decomposed, conditional, dated form in which they are true.

\subsection{References}\label{references}

Cited by bib key against
\texttt{templates/submission/references/pv\_v6.bib} (60 verified entries
as of 2026-07-07). The two entries formerly proposed as UNVERIFIED were
verified and added by A4 on 2026-07-07; the UNVERIFIED markers are
removed at each use:

\begin{itemize}
\tightlist
\item
  \texttt{gopala2008secrecy} (VERIFIED 2026-07-07, in bib): Gopala, P.
  K., Lai, L. and El Gamal, H., ``On the Secrecy Capacity of Fading
  Channels'', IEEE Transactions on Information Theory, 54(10), 2008,
  pp.~4687-4698, doi 10.1109/TIT.2008.928990. Role:
  time-varying-capacity precedent for the time-indexed budgets of ER-6
  and Remark 3.4.
\item
  \texttt{patil2025composition} (VERIFIED 2026-07-07, in bib): Patil,
  V., Stengel-Eskin, E. and Bansal, M., ``The Sum Leaks More Than Its
  Parts: Compositional Privacy Risks and Mitigations in Multi-Agent
  Collaboration'', arXiv:2509.14284. Role: independent
  composition-analysis corroboration of Theorem 4.1's regime, per the
  extraction record.
\end{itemize}

Verified keys used: \texttt{wyner1975wiretap},
\texttt{leungyancheong1978gaussian}, \texttt{csiszar1978broadcast},
\texttt{fano1961transmission}, \texttt{cover2006elements},
\texttt{asif2026infotheoretic}, \texttt{elyagoubi2026agentleak},
\texttt{gopala2008secrecy}, \texttt{patil2025composition},
\texttt{millen1987covert} (section 2), \texttt{sabelfeld2003language}
(section 2), \texttt{dwork2014algorithmic} (section 2).

\end{document}
